% ------------------------------------------------------------------------------
\section{Análisis sísmico}
\label{sec:sis}
% ------------------------------------------------------------------------------

El módulo de análisis sísmico toma la geometría del proyecto activo (las columnas y vigas del modelo 3D), arma con
ella un pórtico tridimensional de concreto armado con un diafragma rígido por piso y lo analiza según la Norma Técnica
E.030 «Diseño Sismorresistente» del Reglamento Nacional de Edificaciones, en su texto 2026 (R.M.~N.\textsuperscript{o}
183-2026-VIVIENDA) o en su versión 2018 (R.M.~N.\textsuperscript{o} 355-2018-VIVIENDA). Este capítulo describe lo que
\textbf{hace el código}, paso a paso, con el nombre de la clase y del método que implementa cada etapa.

\begin{importante}[Estado del módulo]
El análisis sísmico de ConcretoLab es \textbf{educativo y preliminar}. No está validado para diseñar ni para construir
obras, no reemplaza a un programa de análisis estructural completo y sus resultados deben ser revisados por un
ingeniero civil colegiado. La propia aplicación lo advierte en la pantalla, en la lista de advertencias y en el reporte
PDF.
\end{importante}

\begin{nota}[Sobre las citas de la norma]
Los números de artículo y de tabla que aparecen en este capítulo son los que el código devuelve con
\ruta{NormaE030.Parametros.art(...)} o los que escribe en sus mensajes y comentarios. Para esta documentación no se
volvieron a contrastar con el texto oficial; donde hay una duda concreta se indica en el texto.
\end{nota}

% ------------------------------------------------------------------------------
\subsection{Propósito y alcance}
\label{sec:sis-alcance}

El módulo entrega, para un edificio de pórticos de concreto armado:

\begin{itemize}
	\item los parámetros sísmicos $Z$, $U$, $S$, $T_P$, $T_L$, $R_0$, $I_a$, $I_p$, $R$ y $C_T$, elegidos con un asistente de
	      cuatro preguntas;
	\item el peso sísmico, el centro de masas y la inercia polar de cada piso;
	\item todos los modos de vibración del modelo condensado (periodos, formas y masas efectivas);
	\item las fuerzas y cortantes del \textbf{análisis estático} y del \textbf{análisis modal espectral};
	\item los desplazamientos inelásticos, las derivas de entrepiso en el centro de masas y en cada nudo, la razón de
	      torsión $\Delta_{\text{máx}}/\Delta_{\text{prom}}$ y la junta sísmica;
	\item una lista de verificaciones con semáforo y un veredicto («CUMPLE (preliminar)», «NO CUMPLE» o «MODELO
	      INCOMPLETO»);
	\item una recomendación de aumento de las columnas cuando la deriva no cumple (botón «Cómo corregir»);
	\item un análisis tiempo-historia lineal opcional con registros sintéticos o importados;
	\item un reporte PDF con todo lo anterior.
\end{itemize}

Quedan fuera del alcance: muros y placas, sistemas dual, de muros, de muros de ductilidad limitada y de albañilería,
zonas rígidas en los nudos, losas flexibles, interacción suelo-estructura, aislamiento sísmico, análisis no lineal y
diseño de elementos. La figura~\ref{fig:sis-flujo} resume el flujo de cálculo.

\begin{figure}[H]
	\centering
	\begin{tikzpicture}[font=\sffamily\small, caja/.style={draw=clazul, fill=clazulclaro, rounded corners=3pt,
			minimum width=11.5cm, minimum height=0.85cm, align=center}, flecha/.style={-latex, thick, clazuloscuro}]
		\node[caja] (a) at (0,0) {\textbf{Asistente} (\ruta{GuiaSismica}): zona, uso, suelo y sistema $\rightarrow$ \ruta{NormaE030.Parametros}};
		\node[caja] (b) at (0,-1.25) {\textbf{Geometría}: pisos por la cara superior de vigas, columnas y vigas como barras 3D};
		\node[caja] (c) at (0,-2.5) {\textbf{Rigidez} con diafragma rígido y condensación estática; \textbf{pesos y masas} por piso};
		\node[caja] (d) at (0,-3.75) {\textbf{Modos} (Commons Math) $\rightarrow$ análisis \textbf{estático} y \textbf{modal espectral} (CQC)};
		\node[caja] (e) at (0,-5.0) {\textbf{Derivas}, torsión, junta, \ruta{ConexionVigas} $\rightarrow$ verificaciones y veredicto};
		\node[caja, fill=clgris, draw=cltexto] (f) at (0,-6.25) {Recomendación de mejora $\cdot$ Tiempo-historia $\cdot$ Reporte PDF};
		\draw[flecha] (a) -- (b);
		\draw[flecha] (b) -- (c);
		\draw[flecha] (c) -- (d);
		\draw[flecha] (d) -- (e);
		\draw[flecha] (e) -- (f);
	\end{tikzpicture}
	\caption{Flujo del análisis sísmico.}
	\label{fig:sis-flujo}
\end{figure}

% ------------------------------------------------------------------------------
\subsection{Clases que intervienen}
\label{sec:sis-clases}

Todas las clases de cálculo están en \ruta{com.concretolab.service} y no dependen de JavaFX; se prueban sin interfaz
gráfica. La lógica no gráfica del asistente está en \ruta{com.concretolab.controller.GuiaSismica} y la de la pestaña de
tiempo-historia en \ruta{com.concretolab.ui.\allowbreak{}ResumenSimulacion}.

\begin{longtable}{L{6.2cm}L{8.0cm}}
	\caption{Clases del módulo de análisis sísmico.} \label{tab:sis-clases} \\
	\hline
	\textbf{Clase} & \textbf{Responsabilidad} \\ \hline
	\endfirsthead
	\hline
	\textbf{Clase} & \textbf{Responsabilidad} \\ \hline
	\endhead
	\hline
	\endfoot
	\multicolumn{2}{l}{\textit{Norma y parámetros (\ruta{service})}} \\
	\ruta{NormaE030} & Enumeraciones \ruta{Version}, \ruta{Zona}, \ruta{Categoria} y \ruta{Suelo}; registro \ruta{Parametros}
	con $C(T)$, $S_a/g$, $V/P$, exponente $k$, factor inelástico, junta, restricciones de irregularidad y citas de
	artículos (\ruta{art}); límite de deriva \ruta{FERNANDO\_HUANCA\_DERIVA\_LIMITE\_CONCRETO} $= 0{,}007$. \\
	\ruta{TipoSuelo} & Suelo descrito en lenguaje simple y su perfil equivalente en cada versión. \\
	\ruta{UsoEdificacion} & Uso de la edificación: categoría, cargas vivas de la E.020, carga muerta adicional sugerida,
	estado de verificación y si exige aislamiento (A1). \\
	\ruta{SistemaEstructural} & $R_0$ por versión, $C_T$ y si la aplicación puede modelar el sistema (solo pórticos). \\
	\ruta{ZonificacionService} & Bean \ruta{@Service}: zona sísmica por distrito leída de un archivo CSV. \\
	\multicolumn{2}{l}{\textit{Análisis (\ruta{service})}} \\
	\ruta{FernandoHuancaAnalisisSismico} & Motor del análisis: pisos, pórtico 3D, diafragma rígido, condensación, modos,
	estático, espectral, derivas, junta, advertencias, \ruta{ConexionVigas} y \ruta{ModeloDinamico}. \\
	\ruta{RecomendacionSismica} & «Cómo corregir»: menor aumento uniforme de las columnas que hace cumplir la deriva. \\
	\ruta{ReporteSismo} & Reporte PDF del análisis (y del tiempo-historia, si se ejecutó). \\
	\multicolumn{2}{l}{\textit{Tiempo-historia (\ruta{service})}} \\
	\ruta{RegistroSismico} & Acelerograma inmutable de dos componentes con integración y corrección de línea base. \\
	\ruta{GeneradorAcelerograma} & Acelerogramas sintéticos compatibles con el espectro de diseño con $R = 1$. \\
	\ruta{ImportadorRegistro} & Lectura de acelerogramas en archivos de texto (CISMID, IGP, CSMIP, PEER-NGA y similares). \\
	\ruta{EspectroRespuesta} & Espectro de pseudoaceleración de un registro (método de Nair--Jennings). \\
	\ruta{EscaladoRegistros} & Escalado de los conjuntos de registros al espectro de diseño con $R = 1$. \\
	\ruta{AnalisisTiempoHistoria} & Integración de Newmark con amortiguamiento de Rayleigh y verificación de derivas. \\
	\ruta{ResultadoTiempoHistoria} & Historias de un registro: desplazamientos, derivas y cortante basal. \\
	\ruta{VerificacionTiempoHistoria} & Deriva de diseño por piso y dirección, criterio y conjuntos mínimos. \\
	\ruta{ResumenTiempoHistoria} & Resumen de una simulación: escalado, resultados, verificación, tabla y veredicto. \\
	\multicolumn{2}{l}{\textit{Interfaz (\ruta{controller} y \ruta{ui})}} \\
	\ruta{GuiaSismica} & Lógica no gráfica del asistente: parámetros, Tabla 9, cargas por uso, faltantes y verificaciones. \\
	\ruta{WalterQuispesayhuaSismoController} & Pantalla \ruta{sismo-view.fxml}: asistente, cálculo, pestañas de resultados,
	«Cómo corregir» y exportación PDF. \\
	\ruta{TiempoHistoriaController} & Pestaña «Simulación temporal» (\ruta{tiempo-historia-view.fxml}). \\
	\ruta{ResumenSimulacion} & Simulación en segundo plano, periodo para el escalado, decimación de series y colores de
	deriva. \\
	\ruta{VistaModal} & Pestaña «Simulación 3D»: animación de las formas modales sobre el modelo. \\
	\multicolumn{2}{l}{\textit{Auxiliares (sección~\ref{sec:sis-auxiliares})}} \\
	\ruta{DiagnosticoSismico}, \ruta{EdificioCortante}, \ruta{HistoriaPiso}, \ruta{OsciladorSismico}, \ruta{AnalisisPortico}
	& Cálculos simplificados que hoy no usa ninguna pantalla. \\
\end{longtable}

% ------------------------------------------------------------------------------
\subsection{Datos de entrada}
\label{sec:sis-entrada}

\subsubsection{El asistente de cuatro pasos}

La pantalla \ruta{sismo-view.fxml} pide los datos en cuatro tarjetas y un panel plegado de opciones avanzadas
(tabla~\ref{tab:sis-asistente}). \ruta{GuiaSismica.\allowbreak{}walterQuispesayhuaFaltantes} arma la lista de lo que falta: la
ubicación o zona, el uso, el suelo, el sistema (o un sistema que la aplicación pueda modelar) y pisos en el proyecto.
Mientras falte algo, la pantalla lo muestra («Falta: \ldots») y desactiva el botón «Calcular análisis sísmico»; además,
\ruta{onCalcular} marca con error los campos vacíos y no calcula.

\begin{table}[H]
	\centering
	\begin{tabular}{L{3.3cm}L{4.6cm}L{5.6cm}}
		\hline
		\textbf{Paso} & \textbf{Control} & \textbf{Qué produce} \\ \hline
		1. ¿Dónde está la obra? & Departamento, provincia y distrito en cascada, o búsqueda por texto (al menos 2 letras,
		8 resultados), o «zona manual» & Zona sísmica y factor $Z$ (\ruta{ZonificacionService}) \\
		2. ¿Para qué se usará? & Lista de \ruta{UsoEdificacion} & Categoría, $U$, porcentaje de carga viva y cargas por
		piso precargadas \\
		3. ¿Cómo es el suelo? & Lista de \ruta{TipoSuelo} & Perfil de suelo de la versión elegida: $S$, $T_P$, $T_L$ \\
		4. ¿Qué sistema estructural tiene? & Lista de \ruta{SistemaEstructural} & $R_0$ y $C_T$ (solo pórticos se calcula) \\ \hline
	\end{tabular}
	\caption{Pasos del asistente sísmico.}
	\label{tab:sis-asistente}
\end{table}

El panel «Opciones avanzadas (ya tienen valores seguros)» contiene: la versión de la norma (2026 por defecto), $V_{s30}$
(opcional), $I_a \in \{1{,}0;\ 0{,}90;\ 0{,}80;\ 0{,}75;\ 0{,}60;\ 0{,}50\}$,
$I_p \in \{1{,}0;\ 0{,}90;\ 0{,}85;\ 0{,}75;\ 0{,}60\}$, $f'_c = 21$~MPa, coeficiente de Poisson $0{,}20$, fracción de rigidez $EI$ igual a $1{,}0$, espesor de losa
$0{,}20$~m y $C_T$, que se llena solo al elegir el sistema.

Debajo hay una tabla editable con una fila por piso detectado: cota, área de losa (por defecto, el rectángulo envolvente
del piso), carga muerta adicional (CM) y carga viva (CV) en kN/m$^2$. \ruta{GuiaSismica.\allowbreak{}walterQuispesayhuaCargasPorUso}
la precarga al elegir el uso: todos los pisos con la CM adicional sugerida y la CV de pisos, salvo el último, que recibe
la CV de techo. El usuario puede editar cualquier valor.

\subsubsection{De las respuestas a los parámetros}

\ruta{GuiaSismica.\allowbreak{}walterQuispesayhuaParametros} construye el registro \ruta{NormaE030.Parametros} con:

\begin{itemize}
	\item la zona elegida y la categoría del uso (\ruta{uso.categoria()});
	\item el perfil \ruta{tipoSuelo.suelo(version)};
	\item $R_0 = $ \ruta{sistema.r0(version)} y $C_T$ = el valor escrito, o \ruta{sistema.ct()} si es nulo;
	\item $V_{s30}$ solo si la versión es 2026 y el perfil interpola (S2 o S3); en otro caso, \ruta{NaN}.
\end{itemize}

El constructor compacto de \ruta{Parametros} valida que haya versión, zona, categoría y suelo; que el suelo sea de la
misma versión; que $R_0 > 0$, $0 < I_a \le 1$, $0 < I_p \le 1$ y $C_T > 0$; que no se combine el suelo S4 con la zona 4
(la Tabla 4 no tabula $S$ y exige un análisis de respuesta de sitio), y que $V_{s30}$, si se da, esté dentro del intervalo
del perfil. El coeficiente de reducción y la regularidad son
\begin{equation}
	R = R_0 \cdot I_a \cdot I_p, \qquad \text{regular} \iff I_a = I_p = 1 .
	\label{eq:sis-r}
\end{equation}

Las tablas~\ref{tab:sis-zona} a~\ref{tab:sis-uso} reproducen los valores que el código tiene escritos.

\begin{table}[H]
	\centering
	\begin{tabular}{L{2.2cm}L{1.6cm}L{1.0cm}L{6.0cm}}
		\hline
		\textbf{Zona / categoría} & \textbf{Factor} & \textbf{\%CV} & \textbf{Notas} \\ \hline
		Zona 4 & $Z = 0{,}45$ & --- & Tabla 1 \\
		Zona 3 & $Z = 0{,}35$ & --- & \\
		Zona 2 & $Z = 0{,}25$ & --- & \\
		Zona 1 & $Z = 0{,}10$ & --- & \\ \hline
		A (esenciales) & $U = 1{,}5$ & 50 & Tabla 7; la A1 (salud de 2.\textsuperscript{o} y 3.\textsuperscript{er}
		nivel) se representa también con A \\
		B (importantes) & $U = 1{,}3$ & 50 & \\
		C (comunes) & $U = 1{,}0$ & 25 & \\ \hline
	\end{tabular}
	\caption{Zonas (\ruta{NormaE030.Zona}) y categorías (\ruta{NormaE030.Categoria}).}
	\label{tab:sis-zona}
\end{table}

\begin{table}[H]
	\centering
	\begin{tabular}{L{1.3cm}L{1.7cm}L{1.7cm}L{1.7cm}L{1.7cm}L{1.6cm}L{1.6cm}}
		\hline
		\textbf{Perfil} & \textbf{$S$ (Z4)} & \textbf{$S$ (Z3)} & \textbf{$S$ (Z2)} & \textbf{$S$ (Z1)} & \textbf{$T_P$ (s)} &
		\textbf{$T_L$ (s)} \\ \hline
		\multicolumn{7}{l}{\textit{Texto 2026 (intervalo $a$--$b$: $a$ para la mayor velocidad, $b$ para la menor)}} \\
		S0 & 0,80 & 0,80 & 0,80 & 0,80 & 0,3 & 3,0 \\
		S1 & 1,00 & 1,00 & 1,00 & 1,00 & 0,4 & 2,5 \\
		S2 & 1,00--1,10 & 1,00--1,15 & 1,00--1,30 & 1,00--1,30 & 0,4--0,6 & 2,5--2,0 \\
		S3 & 1,10--1,20 & 1,15--1,20 & 1,30--1,40 & 1,30--1,60 & 0,6--0,9 & 2,0--1,6 \\
		S4 & no tabulado & 1,30 & 1,70 & 2,40 & 1,2 & 1,6 \\
		\multicolumn{7}{l}{\textit{Versión 2018 (valores fijos)}} \\
		S0 & 0,80 & 0,80 & 0,80 & 0,80 & 0,3 & 3,0 \\
		S1 & 1,00 & 1,00 & 1,00 & 1,00 & 0,4 & 2,5 \\
		S2 & 1,05 & 1,15 & 1,20 & 1,60 & 0,6 & 2,0 \\
		S3 & 1,10 & 1,20 & 1,40 & 2,00 & 1,0 & 1,6 \\ \hline
	\end{tabular}
	\caption{Perfiles de suelo de \ruta{NormaE030.Suelo} (S2 y S3 de 2026 interpolan con $V_{s30}$ entre 550--350 y
		350--200~m/s).}
	\label{tab:sis-suelo}
\end{table}

\begin{table}[H]
	\centering
	\begin{tabular}{L{3.6cm}L{1.6cm}L{1.6cm}L{6.4cm}}
		\hline
		\textbf{\ruta{TipoSuelo}} & \textbf{2026} & \textbf{2018} & \textbf{Aviso que muestra} \\ \hline
		Roca sana o muy dura & S0 & S1 & En 2018 la roca dura S0 exige $V_{s30} > 1500$~m/s; sin medirlo se toma S1 \\
		Suelo muy firme & S1 & S1 & --- \\
		Suelo firme & S2 & S2 & --- \\
		Suelo intermedio & S3 & S2 & --- \\
		Suelo blando & S4 & S3 & En zona 4 el texto 2026 no tabula $S$ para S4 \\
		No tengo estudio de suelos & S3 & S3 & El art. 16.1 (2026) exige estudio de suelos; resultado referencial \\ \hline
	\end{tabular}
	\caption{Traducción de \ruta{TipoSuelo} al perfil de cada versión.}
	\label{tab:sis-tiposuelo}
\end{table}

\begin{table}[H]
	\centering
	\begin{tabular}{L{4.2cm}L{1.5cm}L{1.5cm}L{0.9cm}L{3.9cm}}
		\hline
		\textbf{\ruta{SistemaEstructural}} & \textbf{$R_0$ 2026} & \textbf{$R_0$ 2018} & \textbf{$C_T$} & \textbf{Modelable} \\ \hline
		Pórticos & 8 & 8 & 35 & Sí (el único) \\
		Dual & 7 & 7 & 60 & No \\
		Muros estructurales & 6 & 6 & 60 & No \\
		Muros de ductilidad limitada & 3,5 & 4 & 60 & No \\
		Albañilería armada o confinada & 3 & 3 & 60 & No \\ \hline
	\end{tabular}
	\caption{Sistemas estructurales (el código cita la Tabla 8 del art. 20, la Tabla 10 del art. 22.1 y el art. 36.1).}
	\label{tab:sis-sistema}
\end{table}

\begin{longtable}{L{4.0cm}L{0.9cm}L{1.3cm}L{1.3cm}L{1.7cm}L{1.9cm}}
	\caption{Usos (\ruta{UsoEdificacion}): categoría y cargas en kN/m$^2$.} \label{tab:sis-uso} \\
	\hline
	\textbf{Uso} & \textbf{Cat.} & \textbf{CV pisos} & \textbf{CV techo} & \textbf{CM adic. sugerida} &
	\textbf{Verificado} \\ \hline
	\endfirsthead
	\hline
	\textbf{Uso} & \textbf{Cat.} & \textbf{CV pisos} & \textbf{CV techo} & \textbf{CM adic.} & \textbf{Verif.} \\ \hline
	\endhead
	\hline
	\endfoot
	Vivienda & C & 2,0 & 1,0 & 2,0 & Sí \\
	Oficinas & C & 2,5 & 1,0 & 1,5 & Sí \\
	Hotel o hostal & C & 2,0 & 1,0 & 2,0 & Sí \\
	Restaurante & C & 4,0 & 1,0 & 1,5 & Sí \\
	Colegio, instituto o universidad & A & 2,5 & 1,0 & 1,5 & Sí \\
	Posta o centro de salud (1.\textsuperscript{er} nivel) & A & 2,0 & 1,0 & 2,0 & Sí \\
	Hospital (2.\textsuperscript{o} y 3.\textsuperscript{er} nivel, A1) & A & 3,0 & 1,0 & 2,0 & Sí \\
	Centro comercial, cine, estadio o iglesia & B & 5,0 & 1,0 & 1,0 & Sí \\
	Industria o taller común & C & 5,0 & 1,0 & 1,0 & No \\
	Depósito o almacén & C & 5,0 & 1,0 & 0,5 & No \\
	Edificación esencial & A & 2,5 & 1,0 & 1,5 & No \\
\end{longtable}

La carga muerta adicional (acabados y tabiquería) \textbf{no proviene de la E.020}: es un valor típico sugerido. Un uso
«no verificado» tiene al menos una carga viva que no figura literalmente en la Tabla 1 de la E.020; el campo
\ruta{fuente} lo explica y la pantalla muestra una nota. Cada uso guarda además la carga viva de corredores y escaleras
(\ruta{cvCorredoresEscaleras}), pero el análisis no la usa. La fuente del uso «Depósito o almacén» menciona el 80~\% del
peso almacenable (art. 31c), pero el cálculo del peso sísmico aplica el porcentaje de su categoría (C, 25~\%); ese 80~\%
no está implementado.

\subsubsection{Versiones 2026 y 2018}

La versión se elige en las opciones avanzadas. Al cambiarla, \ruta{onCambiarVersion} habilita o desactiva la búsqueda
por distrito, la zona manual y $V_{s30}$, y borra los resultados. La tabla~\ref{tab:sis-versiones} lista las diferencias
que implementa el código y la tabla~\ref{tab:sis-articulos}, las citas que devuelve \ruta{Parametros.art}.

\begin{longtable}{L{4.0cm}L{4.9cm}L{4.9cm}}
	\caption{Diferencias entre las versiones en el código.} \label{tab:sis-versiones} \\
	\hline
	\textbf{Aspecto} & \textbf{Texto 2026} & \textbf{Versión 2018} \\ \hline
	\endfirsthead
	\hline
	\textbf{Aspecto} & \textbf{Texto 2026} & \textbf{Versión 2018} \\ \hline
	\endhead
	\hline
	\endfoot
	Perfiles de suelo & S0 a S4; S2 y S3 interpolan con $V_{s30}$ & S0 a S3, valores fijos (los nombres S2 y S3 no
	significan lo mismo que en 2026) \\
	Factor $C$ & Tramo $1 + 7{,}5\,T/T_P$ para $T < 0{,}2\,T_P$ & $C = 2{,}5$ hasta $T_P$ \\
	Combinación direccional & 100~\% en la dirección analizada y 30~\% en la perpendicular & Cada dirección al 100~\%,
	sin componente perpendicular \\
	Zona por distrito & Anexo II, desde el CSV & Desactivada: la zona se elige a mano \\
	Periodo del sitio & Aviso del art. 14.8 en categorías A y B, zona 4, suelos S1 a S4 & Sin aviso \\
	Junta mínima & $0{,}02\,Z S h \ge 0{,}03$~m & $0{,}006\,h \ge 0{,}03$~m \\
	$R_0$ de muros de ductilidad limitada & 3,5 & 4 \\
	Tablas citadas & Tabla 9 (sistemas), Tabla 12 (torsión), Tabla 13 (restricciones), Tabla 14 (deriva) & Tabla 6,
	Tabla 9, Tabla 10, Tabla 11 \\
	Tiempo-historia & Mínimo 7 conjuntos; siempre el promedio & Mínimo 3; promedio con 7 o más, máximo con menos \\
\end{longtable}

\begin{table}[H]
	\centering
	\begin{tabular}{L{4.1cm}L{4.1cm}L{2.4cm}L{2.4cm}}
		\hline
		\textbf{Clave de \ruta{art}} & \textbf{Tema} & \textbf{2026} & \textbf{2018} \\ \hline
		\ruta{aplicabilidad} & Uso del método estático & art. 33.2 & art. 4.5.1 \\
		\ruta{peso} & Peso sísmico & art. 31 & art. 4.3 \\
		\ruta{estatico} & Método estático & arts. 33 a 37 & art. 4.5 \\
		\ruta{excentricidad} & Excentricidad accidental (estático) & art. 37 & art. 4.5.5 \\
		\ruta{espectral} & Análisis modal espectral & arts. 39 a 45 & art. 4.6 \\
		\ruta{cqc} & Combinación modal & art. 42 & art. 4.6.3 \\
		\ruta{minimo} & Cortante mínimo dinámico & art. 44 & art. 4.6.4 \\
		\ruta{excentricidadEspectral} & Excentricidad accidental (espectral) & art. 45 & art. 4.6.5 \\
		\ruta{desplazamientos} & Desplazamientos laterales & art. 50.3 & art. 5.1 \\
		\ruta{junta} & Separación entre edificios & art. 52 & art. 5.3 \\
		\ruta{inerciaBruta} & Secciones brutas & arts. 30.2 y 36.2 & arts. 4.2 y 4.5.4 \\
		\ruta{tablaIrregularidad} & Restricciones a las irregularidades & Tabla 13 & Tabla 10 \\ \hline
	\end{tabular}
	\caption{Citas que devuelve \ruta{NormaE030.Parametros.art(tema)}.}
	\label{tab:sis-articulos}
\end{table}

\subsubsection{Interpolación con \texorpdfstring{$V_{s30}$}{Vs30}}

En el texto 2026, los perfiles S2 y S3 guardan cada valor como un intervalo $\{a, b\}$: $a$ corresponde al extremo de
mayor velocidad ($V_{\text{alto}}$) y $b$ al de menor ($V_{\text{bajo}}$). El método privado \ruta{Suelo.interpolar}
calcula
\begin{equation}
	t = \min\!\left(1,\ \max\!\left(0,\ \frac{V_{\text{alto}} - V_{s30}}{V_{\text{alto}} - V_{\text{bajo}}}\right)\right),
	\qquad X = a + t\,(b - a), \qquad X \in \{S,\ T_P,\ T_L\} .
	\label{eq:sis-vs30}
\end{equation}
Con S2 los extremos son 550 y 350~m/s; con S3, 350 y 200~m/s. Sin $V_{s30}$ se toma $b$, el valor más desfavorable
(art. 17, nota *), y el análisis lo informa en las advertencias. Por ejemplo, S2 en zona 4 con $V_{s30} = 450$~m/s da
$S = 1{,}05$, $T_P = 0{,}5$~s y $T_L = 2{,}25$~s (valores fijados por \ruta{SismoE030Test}). En la versión 2018 el campo
$V_{s30}$ está desactivado y se ignora.

\subsubsection{Zonificación por distrito}

\ruta{ZonificacionService} lee el recurso \ruta{/data/zonificacion-e030-2026.csv}, que en el código fuente está en
\ruta{src/main/resources/data/zonificacion-e030-2026.csv} y viaja dentro del JAR. Es un archivo UTF-8 (se tolera la
marca BOM) con separador punto y coma y cabecera \ruta{departamento;provincia;distrito;zona}; los nombres se conservan
como los escribe el Anexo II del texto 2026, en el mismo orden. Contiene 25 departamentos, 196 provincias y 1892
distritos (cifras que comprueba \ruta{ZonificacionServiceTest}).

\begin{itemize}
	\item La carga es perezosa y única (doble comprobación sobre un campo \ruta{volatile}); los índices quedan inmutables.
	\item Una fila con un número de columnas distinto de 4 o una zona fuera de 1 a 4 lanza \ruta{IllegalStateException}
	      con el número de línea.
	\item Las comparaciones ignoran tildes, mayúsculas y espacios sobrantes (normalización NFD).
	\item \ruta{fernandoHuancaBuscar} ordena los resultados así: distritos que empiezan por el texto, provincias que
	      empiezan por él, distritos que lo contienen y provincias que lo contienen.
	\item \ruta{fernandoHuancaZona} devuelve la zona como \ruta{OptionalInt}; \ruta{GuiaSismica.\allowbreak{}walterQuispesayhuaZonaDe}
	      la convierte en \ruta{NormaE030.Zona}.
\end{itemize}

Con la versión 2018 la búsqueda por distrito se desactiva y la zona manual se fuerza, porque el anexo de esa versión es
otro y la aplicación no lo tiene.

% ------------------------------------------------------------------------------
\subsection{Modelo estructural}
\label{sec:sis-modelo}

El motor es la clase \ruta{FernandoHuancaAnalisisSismico} y su punto de entrada,
\ruta{fernandoHuancaAnalizar(Entrada)}. La \ruta{Entrada} reúne los elementos del proyecto, $f'_c$ (MPa), el coeficiente
de Poisson, la fracción de rigidez, el espesor de losa, las cargas por piso (\ruta{CargaPiso}: área, CM adicional y CV)
y los parámetros de la norma. Antes de calcular se valida que $10 \le f'_c \le 100$~MPa, que la fracción de rigidez
esté en $(0, 1]$, que el coeficiente de Poisson esté en $[0;\ 0{,}5)$, que el espesor no sea negativo y que haya una
carga por piso con área positiva y cargas no negativas.

Solo se usan los elementos de tipo \ruta{COLUMNA} y \ruta{VIGA}. Cada uno es un prisma con centro $(x, y, z)$ y medidas
\ruta{largo} (en X), \ruta{ancho} (en Y) y \ruta{altura} (en Z). Las zapatas no se modelan (base empotrada); los muros y
las losas generan avisos (sección~\ref{sec:sis-advertencias}). Las constantes del motor son la tolerancia
\ruta{FERNANDO\_HUANCA\_TOL} $= 0{,}05$~m, el peso unitario \ruta{FERNANDO\_HUANCA\_PESO\_CONCRETO} $= 24$~kN/m$^3$ y la
gravedad \ruta{FERNANDO\_HUANCA\_G} $= 9{,}80665$~m/s$^2$.

\subsubsection{Detección de pisos}

\ruta{fernandoHuancaPisos} (que la pantalla usa también para llenar la tabla de cargas) detecta los niveles así:

\begin{enumerate}
	\item Calcula la cara superior de cada viga, $z + h/2$, y ordena esos valores.
	\item Los agrupa: un valor abre un grupo nuevo si supera en más de 0,05~m al \textbf{primero} del grupo actual. La cota
	      del piso es el promedio del grupo (\ruta{fernandoHuancaNiveles}).
	\item La base es la menor cota inferior de las columnas, $z_{\text{base}} = \min(z - h/2)$. Se descartan los niveles
	      que no superan $z_{\text{base}} + 0{,}05$~m.
	\item Para cada piso, \ruta{fernandoHuancaEnvolvente} calcula el rectángulo que envuelve las columnas que llegan al
	      nivel y las vigas del nivel; su área es el área de losa por defecto y sus lados $L_x$, $L_y$ se usan para la
	      excentricidad accidental.
\end{enumerate}

Si no hay columnas o vigas, o ninguna viga está por encima de la base, se lanza una \ruta{IllegalArgumentException} con
un mensaje para el usuario. Las alturas de entrepiso son $h_1 = \text{cota}_1 - z_{\text{base}}$ y
$h_i = \text{cota}_i - \text{cota}_{i-1}$, y la altura total es $h_n = \text{cota}_N - z_{\text{base}}$.

\subsubsection{Columnas: nudos y barras}

Para cada columna, el motor busca los pisos cuya cota cumple
$z_{\text{inf}} + 0{,}05 < \text{cota}_i \le z_{\text{sup}} + 0{,}05$. El nudo inicial está en la base (empotrado) si
$|z_{\text{inf}} - z_{\text{base}}| \le 0{,}05$~m; si no, en el piso donde nace la columna. Si la columna nace en una
cota que no es un piso, se considera empotrada y se agrega un aviso. Después crea una barra por cada tramo entre pisos
consecutivos, de modo que una columna que atraviesa varios pisos queda partida en cada nivel. Una columna que no llega
a ningún nivel se ignora con un aviso.

Los nudos se identifican por sus coordenadas redondeadas al milímetro (\ruta{fernandoHuancaNudo}), así que dos
columnas apiladas comparten el nudo del piso. Cada nudo guarda la sección de la columna que lo creó ($h_x$ = largo,
$h_y$ = ancho), que sirve para decidir qué vigas llegan a él.

\subsubsection{Vigas: unión con las columnas}
\label{sec:sis-vigas}

Una viga es de dirección X si su largo es mayor o igual que su ancho; si no, es de dirección Y. Su piso es aquel cuya
cota coincide con la cara superior de la viga con una tolerancia de 0,05~m; si no hay ninguno, se ignora con un aviso.
\ruta{fernandoHuancaEnLinea} decide si el nudo de una columna (centro $(x_n, y_n)$, sección $h_x \times h_y$) está en la
línea de una viga en X de centro $(x_c, y_c)$ y longitud $L$:
\begin{equation}
	|y_n - y_c| \le \frac{h_y}{2} + 0{,}05
	\qquad \text{y} \qquad
	x_c - \frac{L}{2} - \frac{h_x}{2} - 0{,}05 \;\le\; x_n \;\le\; x_c + \frac{L}{2} + \frac{h_x}{2} + 0{,}05 ;
	\label{eq:sis-enlinea}
\end{equation}
para una viga en Y se intercambian los papeles de X e Y. Los nudos que cumplen se ordenan a lo largo de la viga y cada
par consecutivo se une con una barra de eje a eje de columna, sin zonas rígidas; una viga larga que cruza tres columnas
da dos barras. Si hay menos de dos nudos, la viga no aporta rigidez y se agrega el aviso «no une dos columnas». La
sección de la barra es $b \times h$, con $b$ igual al ancho (viga en X) o al largo (viga en Y) y $h$ igual a la altura.

\subsubsection{Diafragma rígido y grados de libertad}

Cada piso tiene tres grados de libertad maestros en su centro de masas: $U_{x,i}$, $U_{y,i}$ y $\theta_i$. Los nudos del
piso se ligan a ellos por la condición de diafragma rígido,
\begin{equation}
	u_x = U_{x,i} - (y - y_{cm,i})\,\theta_i, \qquad u_y = U_{y,i} + (x - x_{cm,i})\,\theta_i, \qquad r_z = \theta_i ,
	\label{eq:sis-diafragma}
\end{equation}
y conservan tres grados de libertad propios ($u_z$, $r_x$, $r_y$). Los nudos de la base no tienen grados de libertad
(empotramiento). La numeración global pone primero los $3N$ grados maestros, piso por piso, y luego los propios de los
nudos. Cada grado local de un nudo se guarda como una lista de índices globales con sus coeficientes
(\ruta{fernandoHuancaGdl} y \ruta{fernandoHuancaCoef}), y el ensamblaje reparte cada término de la rigidez con esos
coeficientes.

\subsubsection{Rigidez de las barras}

\ruta{fernandoHuancaRigidezLocal} arma la matriz de $12 \times 12$ de una barra 3D de Euler--Bernoulli con los términos
$EA/L$ (axial), $GJ/L$ (torsión) y, en cada plano de flexión,
\begin{equation}
	\frac{12EI}{L^3}, \qquad \frac{6EI}{L^2}, \qquad \frac{4EI}{L}, \qquad \frac{2EI}{L},
	\label{eq:sis-flexion}
\end{equation}
con $E = E_c$ (en kPa) y $G = E / [2(1 + \nu)]$. Las dos inercias se multiplican por la fracción de rigidez $f$; el área
y la constante de torsión no. Para una sección $l_2 \times l_3$ se usa $I_z = l_3 l_2^3/12$, $I_y = l_2 l_3^3/12$ y la
constante de torsión de \ruta{fernandoHuancaTorsion}, con $a = \max(l_2, l_3)$ y $b = \min(l_2, l_3)$:
\begin{equation}
	J = a\,b^3 \left[\frac{1}{3} - 0{,}21\,\frac{b}{a}\left(1 - \frac{b^4}{12\,a^4}\right)\right] .
	\label{eq:sis-torsion}
\end{equation}

\begin{table}[H]
	\centering
	\begin{tabular}{L{2.2cm}L{2.1cm}L{2.1cm}L{2.1cm}L{1.5cm}L{1.5cm}}
		\hline
		\textbf{Barra} & \textbf{$e_1$ (eje)} & \textbf{$e_2$} & \textbf{$e_3$} & \textbf{$l_2$} & \textbf{$l_3$} \\ \hline
		Columna & $(0, 0, 1)$ & $(1, 0, 0)$ & $(0, 1, 0)$ & largo & ancho \\
		Viga en X & $(1, 0, 0)$ & $(0, 1, 0)$ & $(0, 0, 1)$ & ancho & altura \\
		Viga en Y & $(0, 1, 0)$ & $(-1, 0, 0)$ & $(0, 0, 1)$ & largo & altura \\ \hline
	\end{tabular}
	\caption{Ejes locales y dimensiones de sección de cada barra.}
	\label{tab:sis-ejes}
\end{table}

Con estos ejes, la rigidez lateral de una columna en X depende del cubo de su largo y en Y del cubo de su ancho; por eso
la recomendación de mejora agranda el largo para corregir X y el ancho para corregir Y. La matriz global de cada barra es
$K_g = T^{\mathsf{T}} K_l\, T$, con $T$ formada por cuatro bloques de la matriz de cosenos directores $[e_1; e_2; e_3]$.

\subsubsection{Condensación estática}

La matriz ensamblada se divide en el bloque maestro ($m$, los $3N$ grados del diafragma) y el bloque propio ($s$). Se
condensa con una factorización LU de Apache Commons Math 3.6.1:
\begin{equation}
	K_c = K_{mm} - K_{ms}\,K_{ss}^{-1}\,K_{sm}, \qquad K_c \leftarrow \tfrac{1}{2}\left(K_c + K_c^{\mathsf{T}}\right).
	\label{eq:sis-condensacion}
\end{equation}
La segunda expresión elimina la asimetría numérica. Si $K_{ss}$ es singular, el mensaje es «El modelo es inestable o
tiene piezas desconectadas\ldots»; si $K_c$ es singular, «La rigidez lateral del modelo es nula: faltan columnas
empotradas o vigas que unan los pórticos». $K_c$ (de $3N \times 3N$, en kN/m, kN y kN$\cdot$m) es la rigidez que usan el
análisis estático, el modal y el tiempo-historia.

\subsubsection{Pesos y masas}

El peso sísmico de cada piso se arma con contribuciones puntuales (registro privado \ruta{Peso}: peso, posición e
inercia propia):

\begin{itemize}
	\item \textbf{Columnas:} cada tramo de longitud $L$ pesa $\gamma_c \cdot \text{largo} \cdot \text{ancho} \cdot L$, con
	      $\gamma_c = 24$~kN/m$^3$, y se reparte mitad al piso de abajo y mitad al de arriba; la mitad que llega a la base
	      no se cuenta. Su inercia propia se toma nula.
	\item \textbf{Vigas:} el peso completo $w = \gamma_c \cdot \text{largo} \cdot \text{ancho} \cdot \text{altura}$ se asigna
	      al piso de la viga, en su centro, con inercia propia $w\,(L^2 + B^2)/12$, donde $L$ y $B$ son su largo y su ancho
	      en planta.
	      También pesan las vigas que no unen columnas, siempre que estén en un piso reconocido.
	\item \textbf{Losa y cargas:} en el centro del rectángulo envolvente,
	      \begin{equation}
	      	W_{\text{losa},i} = A_i \left(e\,\gamma_c + CM_i + \alpha_i\, CV_i\right), \qquad
	      	I_{\text{losa},i} = W_{\text{losa},i}\, \frac{L_x^2 + L_y^2}{12} \cdot \frac{A_i}{L_x L_y},
	      	\label{eq:sis-losa}
	      \end{equation}
	      con $\alpha_i = 0{,}25$ en el último piso (azotea) y, en los demás, el porcentaje de la categoría: 50~\% en A y B,
	      25~\% en C (el código cita el art. 31). El factor $A_i/(L_x L_y)$ escala el rectángulo envolvente al área de losa
	      indicada.
\end{itemize}

No se descuenta la intersección entre vigas, losa y columnas. Con esas contribuciones,
\begin{equation}
	P_i = \sum_k w_k, \qquad x_{cm,i} = \frac{\sum_k w_k x_k}{P_i}, \qquad
	I_i = \sum_k \left[ I_k + w_k \left( (x_k - x_{cm,i})^2 + (y_k - y_{cm,i})^2 \right) \right],
	\label{eq:sis-pesos}
\end{equation}
y las masas del diafragma son $m_{x,i} = m_{y,i} = P_i/g$ (kN$\cdot$s$^2$/m) y $m_{\theta,i} = I_i/g$
(kN$\cdot$s$^2\cdot$m). El peso total es $P = \sum P_i$.

\begin{nota}
\ruta{FERNANDO\_HUANCA\_PESO\_CONCRETO} vale 24~kN/m$^3$ y su comentario lo asocia a los 2400~kgf/m$^3$ de la E.020, que
equivalen a unos 23,5~kN/m$^3$. El código usa el valor redondeado de 24.
\end{nota}

\subsubsection{Módulo de elasticidad}

\ruta{fernandoHuancaModulo} aplica la expresión de la E.060 con $f'_c$ en kgf/cm$^2$ y devuelve MPa:
\begin{equation}
	E_c = 15\,000 \sqrt{\frac{f'_c}{0{,}0980665}} \times 0{,}0980665 \quad [\text{MPa}],
	\label{eq:sis-ec}
\end{equation}
que para $f'_c = 21$~MPa da unos 21\,526~MPa. \ruta{SismoE030Test.\allowbreak{}moduloElasticoSigueLaE060} lo fija en
$21\,538 \pm 60$~MPa (cercano a $4700\sqrt{f'_c}$).

% ------------------------------------------------------------------------------
\subsection{Modos de vibración}
\label{sec:sis-modos}

El problema $K_c\,\phi = \omega^2 M \phi$, con $M$ diagonal, se lleva a forma simétrica estándar:
\begin{equation}
	\tilde{A} = M^{-1/2} K_c M^{-1/2}, \qquad \tilde{A}\,\psi = \lambda\,\psi, \qquad
	\phi = M^{-1/2}\psi, \qquad \omega = \sqrt{\lambda}, \qquad T = \frac{2\pi}{\omega} .
	\label{eq:sis-modos}
\end{equation}
Los autovalores y autovectores los calcula \ruta{EigenDecomposition} de Commons Math. Se ordenan de menor a mayor
autovalor (de mayor a menor periodo) y \textbf{se conservan los $3N$ modos}: no hay truncamiento. Si algún autovalor
no supera $10^{-9}$, el modelo tiene modos de cuerpo rígido y se detiene el cálculo. Como $\psi$ es ortonormal, las
formas $\phi$ quedan normalizadas respecto de la masa ($\phi^{\mathsf{T}} M \phi = 1$), y los factores de participación y
las masas efectivas son
\begin{equation}
	\Gamma_{x,n} = \sum_{i=1}^{N} m_{x,i}\,\phi_{n,(x,i)}, \qquad
	M^{\text{ef}}_{x,n} = \frac{\Gamma_{x,n}^2}{P/g},
	\label{eq:sis-masaefectiva}
\end{equation}
y lo mismo en Y. Como se usan todos los modos, la masa efectiva acumulada es 100~\% en cada dirección
(\ruta{SismoE030Test} lo comprueba con 9 modos para 3 pisos). El periodo fundamental de cada dirección es el del modo con
mayor masa efectiva en esa dirección; es el que se informa como «T fundamental X $\cdot$ Y».

% ------------------------------------------------------------------------------
\subsection{Análisis estático}
\label{sec:sis-estatico}

El método privado \ruta{estatico} aplica los artículos que el código cita como arts. 33 a 37 (2026) o art. 4.5 (2018).

\paragraph{Periodo, factor $C$ y cortante.} El periodo aproximado es el mismo para X e Y:
\begin{equation}
	T = \frac{h_n}{C_T} .
	\label{eq:sis-periodo}
\end{equation}
El factor de amplificación general es (\ruta{Parametros.c}):
\begin{equation}
	C(T) =
	\begin{cases}
		1 + 7{,}5\,T/T_P & T < 0{,}2\,T_P \quad \text{(solo texto 2026)} \\
		2{,}5 & T \le T_P \\
		2{,}5\,T_P/T & T_P < T < T_L \\
		2{,}5\,T_P T_L/T^2 & T \ge T_L
	\end{cases}
	\label{eq:sis-c}
\end{equation}
Para el cortante estático, \ruta{cEstatico} usa $C = 2{,}5$ si $T < T_P$ (el código cita los arts. 18.3 y 34.1), es
decir, el tramo corto no reduce la fuerza estática. El cortante en la base es (\ruta{coeficienteCortante})
\begin{equation}
	V = Z\,U\,S\,\max\!\left(\frac{C}{R},\ 0{,}11\right) P .
	\label{eq:sis-cortante}
\end{equation}

\paragraph{Distribución en altura.} Con $h_i$ medida desde la base,
\begin{equation}
	k =
	\begin{cases}
		1{,}0 & T \le 0{,}5~\text{s} \\
		\min(2{,}0;\ 0{,}75 + 0{,}5\,T) & T > 0{,}5~\text{s}
	\end{cases}
	\qquad
	F_i = V\,\frac{P_i\,h_i^{k}}{\sum_j P_j\,h_j^{k}}, \qquad Q_i = \sum_{j \ge i} F_j .
	\label{eq:sis-distribucion}
\end{equation}

\paragraph{Fuerzas para desplazamientos.} El mínimo $C/R \ge 0{,}11$ no se aplica a los desplazamientos (el código lo
atribuye al art. 50.3). Las fuerzas que se aplican al modelo para calcularlos son
\begin{equation}
	F^{\Delta}_i = \kappa\,F_i, \qquad \kappa = \frac{Z U S\, C/R}{Z U S \max(C/R;\ 0{,}11)} \le 1 .
	\label{eq:sis-reduccion}
\end{equation}

\paragraph{Excentricidad accidental.} Para cada dirección $d$ se resuelven, con la factorización de $K_c$, dos casos con
el momento torsor $M_{t,i} = \pm 0{,}05\,B_i\,F^{\Delta}_i$ aplicado junto a las fuerzas, donde $B_i$ es la dimensión del
rectángulo envolvente perpendicular a la dirección analizada ($L_y$ para X, $L_x$ para Y). Se resuelve además un caso sin
excentricidad para usarlo como acción perpendicular.

\paragraph{Combinación direccional.} En cada punto de control y para cada signo de la excentricidad, la deriva es la del
caso al 100~\% más el 30~\% de la deriva que produce en esa misma dirección la carga perpendicular (sin excentricidad),
sumando valores absolutos; se toma el mayor de los dos signos:
\begin{equation}
	\delta = \max_{\pm e}\left( \left|\delta_{100}^{\pm e}\right| + \mu \left|\delta_{\perp}\right| \right),
	\qquad \mu = \begin{cases} 0{,}30 & \text{texto 2026} \\ 0 & \text{versión 2018} \end{cases}
	\label{eq:sis-combinacion-estatica}
\end{equation}
El término perpendicular solo es distinto de cero cuando la planta acopla ambas direcciones (por ejemplo, con una columna
más rígida que las demás). Las citas del código para esta combinación no coinciden del todo: el comentario del método
\ruta{estatico} menciona los arts. 28.3 y 33.3, mientras que \ruta{fraccionPerpendicular} y el aviso mencionan los
arts. 28.1, 33.3 y 43; verifíquese contra el texto de la norma.

\paragraph{Aplicabilidad.} El código considera aplicable el método estático a estructuras regulares de hasta 30~m y a
cualquier estructura en la zona 1 (art. 33.2 en 2026; art. 4.5.1 en 2018, con la misma regla). Si no es aplicable, lo
avisa y recomienda usar el resultado espectral, pero \textbf{los resultados estáticos se calculan y entran igual en el
veredicto}.

% ------------------------------------------------------------------------------
\subsection{Análisis modal espectral}
\label{sec:sis-espectral}

El método privado \ruta{espectral} aplica los artículos que el código cita como arts. 39 a 45 (2026) o art. 4.6 (2018).

\paragraph{Espectro y respuesta modal.} Para cada modo $n$ y cada dirección del sismo $L$ (X o Y):
\begin{equation}
	\frac{S_{a,n}}{g} = \frac{Z\,U\,C(T_n)\,S}{R}, \qquad
	u_n = \frac{\Gamma_{L,n}\,S_{a,n}}{\omega_n^2}\,\phi_n, \qquad
	f_{n,i} = m_i\,\phi_{n,(L,i)}\,\Gamma_{L,n}\,S_{a,n}, \qquad
	Q_{n,i} = \sum_{j \ge i} f_{n,j} .
	\label{eq:sis-modal}
\end{equation}
El cortante de piso se combina a partir de los cortantes modales (no sumando fuerzas ya combinadas), y la deriva de cada
punto de control se combina a partir de las derivas modales.

\paragraph{Combinación CQC.} Toda respuesta $r$ se combina con
\begin{equation}
	r = \sqrt{\sum_{m}\sum_{n} \rho_{mn}\,r_m\,r_n}, \qquad
	\rho_{mn} = \frac{8\,\xi^2\,(1 + \beta)\,\beta^{3/2}}{(1 - \beta^2)^2 + 4\,\xi^2\,\beta\,(1 + \beta)^2},
	\qquad \beta = \frac{\omega_n}{\omega_m}, \qquad \xi = 0{,}05,
	\label{eq:sis-cqc}
\end{equation}
que es la expresión de Der Kiureghian para amortiguamiento igual en todos los modos (método privado \ruta{cqc}).

\paragraph{Excentricidad accidental.} Con las fuerzas de piso ya combinadas y sin escalar, se aplica en cada piso el
momento $M_{t,i} = 0{,}05\,B_i\,F_i^{\text{CQC}}$, con el mismo signo en todos los pisos, y se resuelve un caso estático.
El valor absoluto de la deriva que produce se suma a la deriva CQC (el código cita el art. 45).

\paragraph{Combinación direccional.} En el texto 2026, la deriva de la dirección $d$ combina por raíz de la suma de
cuadrados el resultado al 100~\% (con su excentricidad) y el 30~\% de la deriva que el espectro perpendicular produce en
esa dirección:
\begin{equation}
	\delta = \sqrt{\left(\delta^{\text{CQC}}_{100} + \left|\delta_{t}\right|\right)^2 + \left(\mu\,\delta^{\text{CQC}}_{\perp}\right)^2} .
	\label{eq:sis-combinacion-espectral}
\end{equation}

\paragraph{Cortante mínimo.} Si el cortante dinámico en la base no alcanza una fracción del estático
(\ruta{fraccionMinimaDinamica}: 80~\% si la estructura es regular, 90~\% si es irregular; art. 44 en 2026, art. 4.6.4 en
2018), las fuerzas y los cortantes se multiplican por
\begin{equation}
	f_{\text{esc}} = \max\!\left(1,\ \frac{\varphi_{\min}\,V_{\text{est}}}{V_{\text{din}}}\right),
	\label{eq:sis-escala}
\end{equation}
donde $V_{\text{est}}$ es el cortante de la ecuación~\eqref{eq:sis-cortante}. \textbf{Los desplazamientos y las derivas
no se escalan}. El periodo que se guarda en el resultado espectral de cada dirección es el del modo con mayor masa
efectiva en ella.

% ------------------------------------------------------------------------------
\subsection{Desplazamientos, derivas, torsión y junta}
\label{sec:sis-derivas}

\paragraph{Desplazamientos inelásticos.} Los desplazamientos elásticos se multiplican por
\ruta{Parametros.\allowbreak{}factorInelastico}:
\begin{equation}
	\Delta_{\text{inel}} = f_{\text{in}}\,\Delta_{\text{el}}, \qquad
	f_{\text{in}} = \begin{cases} 0{,}75\,R & \text{estructura regular} \\ 0{,}85\,R & \text{estructura irregular} \end{cases}
	\label{eq:sis-inelastico}
\end{equation}
El código cita el art. 50.3 (2026) o el art. 5.1 (2018) para los desplazamientos laterales.

\paragraph{Derivas.} Los puntos de control de cada piso son su centro de masas y todos sus nudos. La deriva en un punto
$(x, y)$ es
\begin{equation}
	\delta_i(x, y) = \frac{u_i(x, y) - u_{i-1}(x, y)}{h_i} \le 0{,}007,
	\label{eq:sis-deriva}
\end{equation}
con $u_i(x, y)$ obtenido del diafragma rígido (ecuación~\ref{eq:sis-diafragma}) y $u_0 = 0$. El límite 0,007 es el de
concreto armado (\ruta{FERNANDO\_HUANCA\_DERIVA\_LIMITE\_CONCRETO}; Tabla 14 en 2026 y Tabla 11 en 2018, según el
código). Cada \ruta{Direccion} guarda las derivas en el centro de masas (\ruta{derivasCm}) y el máximo entre todos los
puntos (\ruta{derivasMaximas}); la verificación usa el máximo, y \ruta{Direccion.cumple()} exige
$\max_i \delta_i \le 0{,}007$. El desplazamiento informado por piso (\ruta{desplazamientosCm}) es el del centro de masas.

\paragraph{Razón de torsión.} \ruta{fernandoHuancaRazonTorsion} toma los dos nudos extremos en la coordenada
perpendicular a la dirección analizada (el de menor y el de mayor $y$ para X) y calcula, con las derivas que incluyen la
excentricidad accidental,
\begin{equation}
	r_t = \frac{\max(\delta_a, \delta_b)}{(\delta_a + \delta_b)/2} .
	\label{eq:sis-razon}
\end{equation}
La razón se calcula en todos los pisos, pero solo se tiene en cuenta en los que la deriva supera la mitad del límite
($0{,}0035$). Si $r_t > 1{,}3$ y el $I_p$
elegido es mayor que 0,75 (o $r_t > 1{,}5$ e $I_p > 0{,}60$), las advertencias señalan una posible irregularidad
torsional (Tabla 12 en 2026, Tabla 9 en 2018, según el código).

\paragraph{Junta sísmica.} El registro \ruta{Junta} guarda la altura $h_n$, el mínimo por norma y los dos tercios del
mayor desplazamiento inelástico del techo entre los cuatro análisis (estático y espectral, X e Y):
\begin{equation}
	s_{\min} = \begin{cases} \max(0{,}02\,Z\,S\,h_n;\ 0{,}03~\text{m}) & \text{texto 2026 (art. 52)} \\
		\max(0{,}006\,h_n;\ 0{,}03~\text{m}) & \text{versión 2018 (art. 5.3)} \end{cases}
	\qquad s_{\text{propio}} = \tfrac{2}{3}\,\Delta_{\text{techo,máx}} .
	\label{eq:sis-junta}
\end{equation}
La aplicación solo informa ambos valores: la separación real debe sumar el aporte del edificio vecino y no ser menor que
el mínimo. La fórmula de 2026 es la que tiene el código y la que fija \ruta{SismoE030Test}; verifíquese contra el texto
de la norma.

% ------------------------------------------------------------------------------
\subsection{Verificaciones y veredicto}
\label{sec:sis-verificaciones}

\subsubsection{Vigas sueltas y pisos sin pórtico}

Un modelo cuyas vigas no llegan a las columnas no es un pórtico: sus columnas trabajan como voladizos aislados y, si son
gruesas, pueden dar derivas bajas que no significan nada. Para detectarlo, \ruta{fernandoHuancaConexionVigas} revisa con
el mismo criterio del análisis (ecuación~\ref{eq:sis-enlinea}) cuántas columnas que llegan al piso de cada viga están en
su línea, y devuelve el registro \ruta{ConexionVigas}:

\begin{itemize}
	\item \ruta{vigas}: número total de vigas;
	\item \ruta{idsSueltas} y \ruta{sueltas}: identificadores y nombres de las vigas que no unen al menos dos columnas
	      (incluidas las que no están en un piso reconocido);
	\item \ruta{pisosSinPortico}: pisos (1 es el primero) en los que ninguna viga une dos columnas;
	\item \ruta{formaPorticos()}: verdadero si no hay pisos sin pórtico y no todas las vigas están sueltas.
\end{itemize}

El resultado del análisis lo incluye (\ruta{Resultado.conexion()}) y \ruta{GuiaSismica.\allowbreak{}walterQuispesayhuaModeloValido}
lo consulta. El mismo método se usa en el editor del modelo: antes de guardar una edición,
\ruta{ModeloController.\allowbreak{}yefriLermaVigasQueSeSueltan} compara la conexión antes y después del cambio y, si alguna viga
dejaría de unir dos columnas, pide confirmación («Vigas que quedarían sueltas»).

\subsubsection{Lista de verificaciones}

\ruta{GuiaSismica.\allowbreak{}walterQuispesayhuaVerificaciones} produce la lista del semáforo (tabla~\ref{tab:sis-verificaciones}).
Cada \ruta{Verificacion} tiene un nivel (\ruta{OK}, \ruta{ALERTA}, \ruta{ERROR} o \ruta{INFO}), un texto, un tipo de
corrección (\ruta{NINGUNA}, \ruta{DERIVA}, \ruta{CONSEJO} o \ruta{MODELO}) y un consejo. La pantalla dibuja un botón
«Cómo corregir» en las correcciones \ruta{DERIVA} y \ruta{CONSEJO}, y «Ver en el modelo 3D» en \ruta{MODELO}.

\begin{longtable}{L{3.6cm}L{5.4cm}L{4.8cm}}
	\caption{Verificaciones, en el orden en que aparecen.} \label{tab:sis-verificaciones} \\
	\hline
	\textbf{Verificación} & \textbf{Nivel} & \textbf{Corrección} \\ \hline
	\endfirsthead
	\hline
	\textbf{Verificación} & \textbf{Nivel} & \textbf{Corrección} \\ \hline
	\endhead
	\hline
	\endfoot
	Modelo incompleto (solo si hay vigas sueltas) & \ruta{ERROR} «Modelo incompleto: $n$ de $m$ vigas no unen dos
	columnas\ldots» si no forma pórticos; \ruta{ALERTA} si forma pórticos pero hay alguna suelta & \ruta{MODELO}: abre el
	modelo 3D con esas vigas seleccionadas y el consejo de cómo medirlas \\
	Deriva estática X, estática Y, espectral X, espectral Y & \ruta{OK} si $\le 0{,}007$; \ruta{ERROR} si no; con el
	modelo incompleto, \ruta{ALERTA} «(no válida: modelo incompleto)» & \ruta{DERIVA} en las que fallan: recomendación de
	aumento de columnas \\
	Cortante mínimo & \ruta{OK} sin escalar; \ruta{ALERTA} si se escaló (muestra el factor) & \ruta{CONSEJO}: no es un error
	del modelo \\
	Torsión $\Delta_{\text{máx}}/\Delta_{\text{prom}}$ & \ruta{OK} hasta 1,3; \ruta{ALERTA} sobre 1,3; \ruta{ERROR} sobre
	1,5 (extrema) & \ruta{CONSEJO}: redistribuir rigidez o revisar $I_p$ \\
	Método estático aplicable & \ruta{OK} o \ruta{ALERTA} & \ruta{CONSEJO}: usar el espectral \\
	Sistema estructural (Tabla 9 en 2026, Tabla 6 en 2018) & \ruta{ERROR} en categoría A y zonas 2 a 4; \ruta{OK} en los
	demás casos & \ruta{CONSEJO}: sistema dual o de muros \\
	Aislamiento (uso A1 en zonas 3 y 4) & \ruta{ERROR} (solo si aplica) & \ruta{CONSEJO}: requiere aisladores \\
	Junta sísmica & \ruta{INFO} con el mínimo y los 2/3 del desplazamiento propio & --- \\
\end{longtable}

La línea de la Tabla 9 llama a \ruta{walterQuispesayhuaAvisoTabla9} sin la marca de uso A1. Por eso, para un hospital
(A1) en zona 1 la lista muestra «los pórticos están permitidos», mientras que el aviso rojo del resumen, que sí recibe
esa marca, dice que no lo están. En las zonas 2 a 4 ambos coinciden, porque toda la categoría A queda prohibida.

La torsión de esta lista usa la mayor razón de los cuatro análisis en los pisos con deriva mayor que $0{,}0035$
(\ruta{walterQuispesayhuaRazonTorsionMaxima}) y no tiene en cuenta el $I_p$ elegido; el aviso de la lista de advertencias,
en cambio, solo aparece si ese $I_p$ todavía no refleja la irregularidad.

\subsubsection{Veredicto}

\ruta{WalterQuispesayhuaSismoController.\allowbreak{}walterQuispesayhuaVeredicto} decide el texto del encabezado del resumen:

\begin{enumerate}
	\item «MODELO INCOMPLETO: las vigas no forman pórticos», si \ruta{walterQuispesayhuaModeloValido} es falso;
	\item «CUMPLE (preliminar)», si las cuatro direcciones (estática X e Y, espectral X e Y) tienen deriva máxima
	      $\le 0{,}007$;
	\item «NO CUMPLE», en otro caso.
\end{enumerate}

El veredicto solo depende de las derivas y de la conexión de las vigas. Una línea en rojo por la Tabla 9, por torsión
extrema o por aislamiento no lo cambia; por eso la lista de verificaciones se muestra siempre debajo. Con el modelo
incompleto, la pestaña «Derivas» y el reporte PDF muestran «No válida» en lugar de «Cumple» o «No cumple».

\subsubsection{Advertencias del análisis}
\label{sec:sis-advertencias}

Además de la lista de verificaciones, \ruta{fernandoHuancaAnalizar} devuelve una lista de textos
(\ruta{Resultado.\allowbreak{}advertencias}) que la pantalla muestra bajo el resumen y el PDF reproduce. Según el caso, incluye:

\begin{itemize}
	\item columnas o vigas ignoradas o consideradas empotradas por estar fuera de los pisos reconocidos, y vigas que no
	      unen dos columnas;
	\item método estático no aplicable y cortante dinámico escalado;
	\item restricciones a las irregularidades (\ruta{Parametros.\allowbreak{}restricciones}, Tabla 13 en 2026 o Tabla 10 en 2018): en
	      categoría A y zonas 2 a 4 no se permiten irregularidades; en A no se permiten extremas; en B y zonas 2 a 4 no se
	      permiten extremas; en C y zonas 3 y 4 no se permiten extremas; en C y zona 2 las extremas solo hasta 2 pisos u
	      8~m. El código considera extrema una irregularidad con $I_a$ o $I_p \le 0{,}60$;
	\item muros en el modelo, que no se incluyen en la rigidez ni en el cortante;
	\item losas del modelo, que solo se usan para el metrado (el peso sale de las cargas por m$^2$);
	\item en 2026, zona 4, categorías A o B y suelos S1 a S4: el art. 14.8 pide verificar que el periodo del sitio sea
	      menor que $0{,}65\,T_P$, cosa que el módulo no calcula;
	\item $V_{s30}$ no indicado en S2 o S3 (con los valores usados);
	\item la combinación direccional aplicada (2026) o la nota de que la versión 2018 no la pide;
	\item posible irregularidad torsional, según la razón $r_t$ y el $I_p$ elegido;
	\item masa efectiva acumulada menor que 90~\% (con todos los modos conservados no debería aparecer);
	\item periodo $h_n/C_T$ distinto en más de 50~\% del periodo modal en X;
	\item rigidez reducida por debajo de 1,0, porque la norma pide secciones brutas;
	\item la junta sísmica y el recordatorio de base empotrada, sin zonas rígidas ni muros.
\end{itemize}

% ------------------------------------------------------------------------------
\subsection{Recomendación de mejora}
\label{sec:sis-recomendacion}

Cuando una deriva no cumple, su línea ofrece «Cómo corregir». El controlador llama una sola vez por resultado a
\ruta{RecomendacionSismica.\allowbreak{}fernandoHuancaRecomendar(Entrada, Resultado)}, que busca el \textbf{menor aumento uniforme} de
la sección de todas las columnas que hace cumplir la deriva:

\begin{enumerate}
	\item Toma la deriva de cada dirección como el mayor valor entre el estático y el espectral
	      (\ruta{fernandoHuancaDerivaX} y \ruta{fernandoHuancaDerivaY}). Si ambas cumplen, responde que no hace falta
	      cambiar nada. Si no hay columnas, sugiere agregar columnas o muros.
	\item Repite: suma \ruta{FERNANDO\_HUANCA\_PASO} $= 0{,}05$~m al largo (X) de todas las columnas si falla X y al ancho
	      (Y) si falla Y, y vuelve a ejecutar \ruta{fernandoHuancaAnalizar} sobre \textbf{copias} de los elementos, con las
	      mismas cargas, materiales y parámetros. Así la medida sale del cálculo e incluye el peso de las columnas más
	      grandes. Las medidas se redondean al milímetro.
	\item Se detiene cuando ambas direcciones cumplen o cuando un aumento superaría \ruta{FERNANDO\_HUANCA\_MAXIMO}
	      $= 0{,}60$~m. En ese caso limita el aumento a 0,60~m, marca la recomendación como insuficiente y sugiere placas o
	      un sistema dual validado en un programa de análisis.
	\item Si es suficiente pero el aumento supera \ruta{FERNANDO\_HUANCA\_AUMENTO\_GRANDE} $= 0{,}30$~m, el texto agrega:
	      «Es un aumento grande: probablemente convenga más peraltar las vigas o agregar placas en esa dirección que
	      agrandar tanto las columnas».
\end{enumerate}

El resultado (\ruta{Recomendacion}) lleva los aumentos en X e Y, la lista de \ruta{Cambio} (columna, largo nuevo y ancho
nuevo), las derivas antes y después, y un texto con las secciones actuales y nuevas (hasta cuatro distintas). El diálogo
«Cómo corregir la deriva» ofrece «Ir al modelo 3D y precargar el cambio»: navega a \ruta{modelo-view} y llama a
\ruta{ModeloController.\allowbreak{}yefriLermaPrecargarCorreccion}. El cambio no se guarda solo: el usuario lo revisa, pulsa
«Aplicar» y vuelve a calcular el análisis.

\begin{consejo}
La recomendación agranda \textbf{todas} las columnas por igual y solo en la dirección que falla. Es una medida de
anteproyecto; otras soluciones (peraltar vigas, agregar placas, redistribuir columnas) pueden ser mejores y la
aplicación lo dice en el mismo texto.
\end{consejo}

% ------------------------------------------------------------------------------
\subsection{Análisis tiempo-historia}
\label{sec:sis-th}

La pestaña «Simulación temporal» está incluida en la vista sísmica (\ruta{fx:include} de \ruta{tiempo-historia-view.fxml})
y se configura después de cada cálculo con \ruta{TiempoHistoriaController.\allowbreak{}walterQuispesayhuaConfigurar}. Usa el
\ruta{ModeloDinamico} que devuelve el análisis: masas, rigidez condensada $K_c$, cotas, alturas, centros de masas,
puntos de control y periodos. Toda la lógica de cálculo está en las clases de servicio; la pantalla solo la orquesta.

\subsubsection{Flujo en la pantalla}

\begin{enumerate}
	\item El periodo para el escalado es $T = \max(T_x, T_y)$ (\ruta{ResumenSimulacion.\allowbreak{}walterQuispesayhuaPeriodoFundamental}),
	      con los periodos fundamentales del análisis espectral; así la rejilla $0{,}2\,T$ a $1{,}5\,T$ cubre ambas
	      direcciones.
	\item El usuario elige una fuente: «Generar» crea 7 conjuntos sintéticos (constante
	      \ruta{WALTER\_QUISPESAYHUA\_CONJUNTOS\_SINTETICOS}) con la semilla indicada (2026 por defecto); o «Importar
	      acelerograma\ldots» agrega registros a una lista (un archivo de tres columnas o dos archivos, uno por componente,
	      con la unidad y, si hace falta, el paso de tiempo) y «Simular importados» los procesa.
	\item Un \ruta{Task} de JavaFX ejecuta en un hilo aparte \ruta{ResumenSimulacion.\allowbreak{}walterQuispesayhuaSimular}: escala los
	      registros, resuelve cada uno con 5~\% de amortiguamiento y verifica las derivas, informando el avance.
	\item La pestaña muestra el veredicto, las advertencias, una tabla por registro (PGA, factor, desplazamiento máximo del
	      techo, deriva máxima y cortante basal máximo en X e Y, más la fila del promedio o máximo y la del límite), la
	      gráfica del espectro de diseño con $R = 1$ frente al promedio SRSS escalado y una animación del edificio que se
	      mueve con el terreno (desplazamiento absoluto = relativo + terreno). Los entrepisos se pintan en verde hasta el
	      60~\% del límite, en ámbar hasta el 100~\% y en rojo por encima.
\end{enumerate}

\subsubsection{Registros sísmicos}

\ruta{RegistroSismico} es un registro inmutable con nombre, fuente, $\Delta t$ (entre 0,001 y 0,1~s) y dos componentes
horizontales de igual longitud (de 2 a 200\,000 pasos) en m/s$^2$. Las velocidades y desplazamientos del terreno se
obtienen con una corrección de línea base (\ruta{integrar}): se resta la aceleración media, se integra la velocidad por
trapecios y se le resta un polinomio $\alpha t + \beta t^2$ elegido para que la velocidad final y el desplazamiento final
sean nulos; luego se integra el desplazamiento. \ruta{escalado(factor)} devuelve una copia con ambas componentes
multiplicadas por el mismo factor.

\paragraph{Registros sintéticos.} \ruta{GeneradorAcelerograma.\allowbreak{}fernandoHuancaGenerarConjuntos} genera, por componente:

\begin{itemize}
	\item una densidad espectral objetivo
	      $G(\omega) = \dfrac{4\zeta}{\pi\omega}\left(\dfrac{S_a(\omega)}{\eta}\right)^2$ con $\zeta = 0{,}05$ y el factor
	      de pico $\eta = \max\!\left(2,\ \sqrt{2 \ln \max(2 f t_s;\ 3)}\right)$, $t_s = 15$~s;
	\item una suma de armónicos con fases aleatorias (FFT inversa de 4096 puntos), $\Delta t = 0{,}01$~s y 40~s de
	      duración (4001 pasos);
	\item una envolvente de intensidad: $(t/3)^2$ hasta 3~s, 1 hasta 15~s y $e^{-0{,}2(t - 15)}$ después;
	\item hasta 25 iteraciones de corrección en frecuencia, que multiplican el espectro de Fourier por el cociente objetivo /
	      espectro de respuesta, hasta que cada componente quede entre 95 y 110~\% del objetivo en $0{,}2\,T$ a $1{,}5\,T$
	      (el criterio interno es 95,5--109,5~\% apuntando a 102,5~\%); se eliminan las frecuencias menores de 0,05~Hz.
\end{itemize}

Las componentes X e Y usan fases independientes y el resultado es reproducible con la misma semilla. Si una componente no
converge, la razón mínima alcanzada se anota en la fuente del registro.

\paragraph{Registros importados.} \ruta{ImportadorRegistro} lee archivos de texto de hasta 10~MB: una columna
(aceleración, con $\Delta t$ indicado o leído del encabezado, por ejemplo «DT=0.005», «DELTA T» o «INTERVALO»), dos
columnas (tiempo y aceleración), tres columnas (tiempo, $a_x$, $a_y$) o más de tres columnas como una sola serie (formato
PEER-NGA, si se conoce $\Delta t$). Admite espacio, tabulación, coma o punto y coma como separadores y coma decimal;
omite el encabezado y sugiere la unidad si el encabezado la menciona. Las unidades son m/s$^2$, cm/s$^2$ (gal, factor
0,01) y $g$ (factor 9,80665). El $\Delta t$ de una columna de tiempo es la mediana de los incrementos y debe ser uniforme
con $\pm 1$~\%. Al importar un par de archivos, ambos deben tener el mismo $\Delta t$ (1~\%) y casi la misma longitud (se
recorta al menor).

\subsubsection{Espectro de respuesta y escalado}

\ruta{EspectroRespuesta.\allowbreak{}fernandoHuancaPseudoAceleracion} resuelve un oscilador amortiguado con el método exacto para
excitación lineal por tramos (Nair--Jennings, coeficientes de la Tabla 5.2.1 de Chopra) y devuelve
$S_a = \omega^2 \max|u(t)|$. Para $T \le 0{,}02$~s toma la aceleración pico del terreno, y si el paso del registro es
grueso respecto del periodo, subdivide cada intervalo (hasta 64 veces) para trabajar con unos $T/20$.

\ruta{EscaladoRegistros.\allowbreak{}fernandoHuancaEscalar} trabaja en una rejilla logarítmica de 60 periodos entre $0{,}2\,T$ y
$1{,}5\,T$ (el código cita el art. 47.5) y con el espectro de diseño del art. 41 con $R = 1$
(\ruta{fernandoHuancaConR1} pone $R_0 = I_a = I_p = 1$). Para cada conjunto calcula el espectro de cada componente al
5~\% y su combinación $\text{SRSS}_i(T) = \sqrt{S_{a,x}^2 + S_{a,y}^2}$ (art. 47.4), y luego:
\begin{align}
	f_i &= \left[\frac{1}{m}\sum_{j=1}^{m} \frac{\text{SRSS}_i(T_j)}{S_{a,d}(T_j)}\right]^{-1}
	&& \text{(registros reales)} \label{eq:sis-fi} \\
	f_i &= \max\!\left(1,\ \frac{0{,}90}{\min_j \min(S_{a,x}, S_{a,y})/S_{a,d}}\right)
	&& \text{(sintéticos compatibles)} \label{eq:sis-fi-compatible} \\
	k_{\text{c}} &= \max\!\left(1,\ \left[\min_j \frac{\frac{1}{n}\sum_i f_i\,\text{SRSS}_i(T_j)}{S_{a,d}(T_j)}\right]^{-1}\right),
	&& \text{factor final} = f_i\,k_{\text{c}} . \label{eq:sis-k}
\end{align}
Aquí $S_{a,d}$ es el espectro de diseño con $R = 1$, $m = 60$ periodos y $n$ conjuntos. El mismo factor se aplica a las
dos componentes de un conjunto. Con el factor común $k_{\text{c}}$, el promedio SRSS escalado no baja del 100~\% del
espectro de diseño en ningún periodo de la rejilla. El escalado avisa si alguna componente escalada queda por debajo del
90~\% del espectro (art. 47.6; con registros reales solo es un aviso), si hay menos conjuntos que el mínimo (7 en el
texto 2026, art. 47.1; 3 en 2018, art. 4.7) y si algún factor es mayor que 4 o menor que 0,25. Cumple si la razón
mínima es al menos 0,999, hay conjuntos suficientes y, con sintéticos, ninguna componente queda bajo el 90~\%.

\subsubsection{Integración en el tiempo}

\ruta{AnalisisTiempoHistoria.\allowbreak{}fernandoHuancaResolver} integra el sistema acoplado de $3N$ grados de libertad con las dos
componentes del registro actuando a la vez:
\begin{equation}
	M\ddot{u} + C\dot{u} + K_c\,u = -M\left(r_x\,a_{g,x}(t) + r_y\,a_{g,y}(t)\right),
	\label{eq:sis-movimiento}
\end{equation}
donde $u$ es el desplazamiento relativo a la base y $r_x$, $r_y$ son los vectores de influencia
(\ruta{ModeloDinamico.influencia}). El amortiguamiento es de Rayleigh, $C = \alpha M + \beta K_c$, calibrado para que
$\xi = 0{,}05$ (art. 48.5) se cumpla en $T_1$ = \ruta{periodos()[0]} (el mayor periodo del modelo) y en $T_2 = 0{,}2\,T_1$;
si $0{,}2\,T_1$ es menor que el menor periodo del modelo, $T_2$ se toma igual a ese periodo:
\begin{equation}
	\alpha = \frac{2\,\xi\,\omega_1\omega_2}{\omega_1 + \omega_2}, \qquad \beta = \frac{2\,\xi}{\omega_1 + \omega_2} .
	\label{eq:sis-rayleigh}
\end{equation}
Ese $T_1$ puede no coincidir con el $T$ del escalado, que es el mayor de los periodos fundamentales en X e Y (por
ejemplo, si el primer modo es torsional). El método es Newmark de aceleración media ($\beta_N = 1/4$, $\gamma_N = 1/2$), incondicionalmente estable. Como
$\Delta t$ es constante, la matriz efectiva se factoriza con LU una sola vez:
\begin{equation}
	\hat{K} = K_c + \frac{2}{\Delta t}\,C + \frac{4}{\Delta t^2}\,M .
	\label{eq:sis-newmark}
\end{equation}
Se parte del reposo, con $\ddot{u}_0 = -r\,a_g(0)$. Si $\Delta t$ supera la décima parte del menor periodo, se avisa de
la pérdida de precisión en los modos altos. El cortante basal de cada instante es la suma de las fuerzas de inercia,
$V_b(t) = -\sum_i m_i\,(\ddot{u}_i + a_g)$. Las derivas se calculan en los mismos puntos de control del análisis
espectral y son \textbf{elásticas}: no se multiplican por $0{,}75\,R$, porque el registro ya está escalado al espectro con
$R = 1$.

\subsubsection{Verificación del artículo 49}

\ruta{AnalisisTiempoHistoria.\allowbreak{}fernandoHuancaVerificar} toma, por piso y dirección, el máximo en el tiempo de cada registro
y los combina:

\begin{itemize}
	\item \textbf{texto 2026}: siempre el promedio de los máximos (art. 49.1), con un mínimo de 7 conjuntos (art. 47.1);
	\item \textbf{versión 2018}: el promedio si hay 7 conjuntos o más y el máximo si hay menos, con un mínimo de 3
	      (art. 4.7.3).
\end{itemize}

El límite es $1{,}25 \times 0{,}007 = 0{,}00875$ (art. 49.2 en 2026, art. 4.7.3 en 2018, según el código). La
verificación cumple si todos los pisos cumplen en X e Y y hay conjuntos suficientes. \ruta{ResumenTiempoHistoria.veredicto}
devuelve «CUMPLE», «NO CUMPLE» o «NO CONCLUYENTE» (menos conjuntos que el mínimo), y el controlador lo reemplaza por
«MODELO INCOMPLETO» si las vigas no forman pórticos. Las advertencias recuerdan siempre que el tiempo-historia es una
verificación complementaria que no sustituye al análisis estático ni al espectral (arts. 46.1 y 32.3), que el modelo es
lineal-elástico con $R = 1$ y que la base está empotrada.

% ------------------------------------------------------------------------------
\subsection{Reporte PDF}
\label{sec:sis-pdf}

El botón «Exportar PDF» llama a \ruta{ReporteSismo.\allowbreak{}fernandoHuancaGenerar(proyecto, entrada, resultado, resumen)}, que
arma el documento con OpenPDF 1.3.43 (paquete \ruta{com.lowagie}) en tamaño A4, y lo guarda con
\ruta{Exportador.\allowbreak{}fernandoHuancaGuardar} (nombre sugerido \ruta{analisis-sismico}). Si se ejecutó una simulación
temporal, se incluye su último resumen. El contenido es:

\begin{enumerate}
	\item título, versión de la norma, proyecto y ubicación;
	\item \textbf{datos y parámetros}: zona, categoría, suelo ($S$, $V_{s30}$ o la nota del valor más desfavorable), $T_P$,
	      $T_L$, $R_0$, $I_a$, $I_p$, $R$, regularidad, $C_T$, límite de deriva, $f'_c$, $E_c$, Poisson, fracción de
	      rigidez, espesor de losa y peso sísmico total;
	\item \textbf{pisos y pesos}: cota, altura, área, peso, centro de masas e inercia polar, con la nota del porcentaje de
	      carga viva y del peso unitario;
	\item \textbf{modos}: periodo, frecuencia y masas efectivas, con la masa acumulada;
	\item \textbf{análisis estático}: $T$, $C$, $k$, $V$, excentricidad y combinación direccional, y la tabla de fuerzas y
	      cortantes por piso;
	\item \textbf{análisis modal espectral}: espectro, CQC, excentricidad, combinación direccional, factores de escala y la
	      tabla de fuerzas y cortantes;
	\item \textbf{desplazamientos y derivas} inelásticos por piso, con estado «Cumple», «No cumple» o «No válida» y, si el
	      modelo está incompleto, un párrafo «MODELO INCOMPLETO»;
	\item \textbf{tiempo-historia} (si existe): origen de los registros, conjuntos, $T$ y rejilla, razón de escalado, tabla
	      por registro, factores, resultado del art. 49, deriva de diseño por piso y advertencias;
	\item \textbf{advertencias} del análisis, \textbf{referencias} (E.030 de la versión elegida, E.060, E.020,
	      Der Kiureghian 1981, McGuire et al. 2000 y Chopra) y una nota final de que los resultados son educativos y
	      preliminares.
\end{enumerate}

% ------------------------------------------------------------------------------
\subsection{Otras clases del paquete}
\label{sec:sis-auxiliares}

El paquete \ruta{service} contiene cinco clases sísmicas simplificadas que \textbf{ninguna pantalla usa hoy}; solo las
llaman sus pruebas (\ruta{AnalisisPorticoTest}, \ruta{HistoriaPisoTest}, \ruta{OsciladorSismicoTest} y
\ruta{ProyectoIntegracionTest}):

\begin{itemize}
	\item \ruta{DiagnosticoSismico.\allowbreak{}yefriLermaRevisar}: cuenta zapatas, columnas y vigas y devuelve una lista fija de lo que
	      faltaría para un modelo analítico. No calcula nada.
	\item \ruta{OsciladorSismico.\allowbreak{}yefriLermaCalcular}: oscilador de un grado de libertad; $T = 2\pi\sqrt{m/k}$,
	      $F = m\,S_a\,g$ y $\Delta = F/k$ (en mm).
	\item \ruta{HistoriaPiso}: rigidez de un piso como suma de $12EI/h^3$ de columnas biempotradas
	      (\ruta{yefriLermaPiso}) e integración de Newmark de aceleración media para un grado de libertad
	      (\ruta{yefriLermaResolver}, exige $\Delta t \le T/20$).
	\item \ruta{EdificioCortante.\allowbreak{}yefriLermaSimular}: edificio de cortante con masas iguales y una traslación por planta
	      (1 a 12 niveles), diagonalización de Jacobi y superposición modal con \ruta{HistoriaPiso}.
	\item \ruta{AnalisisPortico.\allowbreak{}yefriLermaResolver}: pórtico plano elástico por rigidez directa (hasta 60 nudos y 180
	      barras) con eliminación de Gauss y pivoteo parcial; devuelve desplazamientos y reacciones.
\end{itemize}

% ------------------------------------------------------------------------------
\subsection{Pruebas automáticas}
\label{sec:sis-pruebas}

Las pruebas no validan el módulo frente a la norma ni frente a otro programa: comprueban que el código calcula lo que
dice calcular, con valores conocidos y casos con solución analítica. La tabla~\ref{tab:sis-pruebas} resume las seis
clases de prueba del módulo.

\begin{longtable}{L{4.8cm}L{9.6cm}}
	\caption{Pruebas del módulo sísmico y valores que fijan.} \label{tab:sis-pruebas} \\
	\hline
	\textbf{Clase de prueba} & \textbf{Qué comprueba} \\ \hline
	\endfirsthead
	\hline
	\textbf{Clase de prueba} & \textbf{Qué comprueba} \\ \hline
	\endhead
	\hline
	\endfoot
	\ruta{SismoE030Test} & Tablas: S2 en zona 4 sin $V_{s30}$ da $S = 1{,}10$; $C(0{,}4) = 2{,}5$, $C(1{,}0) = 1{,}5$ y
	$C(3{,}0) = 2{,}5 \cdot 0{,}6 \cdot 2{,}0/9$ con $T_P = 0{,}6$ y $T_L = 2{,}0$; $k(0{,}5) = 1$, $k(1{,}0) = 1{,}25$,
	$k(3{,}0) = 2$; factor inelástico 6 ($0{,}75 \times 8$) y 5,1 ($0{,}85 \times 6$). Mínimo $C/R = 0{,}11$. Tramo corto:
	$C(0{,}06) = 1{,}75$ y $C_{\text{est}}(0{,}06) = 2{,}5$. Interpolación: $V_{s30} = 450$ da $S = 1{,}05$,
	$T_P = 0{,}5$, $T_L = 2{,}25$; S3 en zona 1 con 275~m/s da $S = 1{,}45$, $T_P = 0{,}75$, $T_L = 1{,}8$; S4 en zona 4
	y $V_{s30}$ fuera de rango se rechazan. Un piso con vigas casi rígidas tiene un periodo a menos de 12~\% del de cuatro
	columnas biempotradas. Tres pisos: 9 modos ordenados y 100~\% de masa efectiva; $T = 9/35$; las fuerzas estáticas
	suman $V$ y crecen con la altura. El espectral cumple el 80~\% del estático y es simétrico en planta simétrica. Las
	derivas estáticas no dependen de $R$ cuando rige el mínimo 0,11. La excentricidad accidental eleva la deriva máxima
	sobre la del centro de masas. Versión 2018: sus tablas, sin el 30~\%, citas \ruta{art. 4.6.4} y \ruta{art. 44}. Junta:
	$0{,}02 \times 0{,}45 \times 1{,}10 \times 9$ (2026) y $0{,}006 \times 9$ (2018), mínimo 0,03~m. $E_c(21) = 21\,538
	\pm 60$~MPa. El PDF empieza con \ruta{\%PDF}. \\
	\ruta{GuiaSismicaTest} & Tabla 9: pórticos prohibidos en categoría A y zonas 2 a 4, y en A1 en todas las zonas.
	Conversión de zonas. Armado de parámetros (vivienda y suelo intermedio en 2026 dan C, S3, $R_0 = 8$, $C_T = 35$; en
	2018 el suelo intermedio pasa a S2 y $V_{s30}$ se ignora). Cargas por uso con la CV de techo en el último piso. Lista de
	faltantes. La lista de verificaciones incluye la junta (\ruta{INFO}), la Tabla 9 (\ruta{ERROR} en un colegio en zona 4)
	y cuatro líneas de deriva. \\
	\ruta{RecomendacionSismicaTest} & Pórtico de 4 pisos con columnas de $0{,}80 \times 0{,}40$~m que solo falla en Y: la
	recomendación es $+0{,}10$~m en Y y nada en X (con $+0{,}05$ la deriva aún es 0,0075; con $+0{,}10$, 0,0069), cambia
	las 4 columnas, es el mínimo y no modifica los elementos originales. Sin fallas no hay cambios. Las derivas que fallan
	llevan la corrección \ruta{DERIVA}. \\
	\ruta{TiempoHistoriaTest} & Oscilador armónico: amplitud estacionaria con error menor que 1,5~\% y cortante basal menor
	que 2~\%. Sin carga no hay movimiento. El periodo de vibración libre es $T_1$. Planta simétrica sin movimiento en Y ni
	giro. Convergencia entre $\Delta t$ y $\Delta t/2$ dentro de 1~\%. Desplazamiento absoluto = relativo + terreno.
	Verificación: límite 0,00875; 7 conjuntos promedian; en 2018 con 3 conjuntos se toma el máximo; 3 conjuntos en 2026 y
	2 en 2018 son insuficientes. \\
	\ruta{RegistrosSismicosTest} & Validaciones del registro, integración y corrección de línea base. Espectro de un seno
	en resonancia igual a $1/(2\xi)$ con error menor que 1~\%, tendencia a la PGA en periodos cortos e impulso con solución
	analítica. Importación de una, dos, tres columnas y formato PEER, pares de archivos, errores claros, $\Delta t$ no
	uniforme y archivos mayores de 10~MB. Generador reproducible con la semilla, 7 conjuntos de 40~s con
	$\Delta t = 0{,}01$~s y espectro entre 85 y 125~\% del objetivo. Escalado con promedio SRSS $\ge$ diseño, mismo factor
	en ambas componentes y avisos de conjuntos y del art. 47.6. \\
	\ruta{ModeloIncompletoTest} (paquete \ruta{controller}) & Pórtico de 3 pisos (3 $\times$ 2 columnas a 5 $\times$ 4~m) bien armado: sin
	vigas sueltas. Con columnas de 1,1~m y vigas convertidas en bloques de $0{,}6 \times 0{,}6$~m, las cuatro derivas
	«cumplen», pero todas las vigas están sueltas, los pisos 1, 2 y 3 no tienen pórtico, el veredicto empieza con «MODELO
	INCOMPLETO» y las derivas salen en \ruta{ALERTA} «no válida». Una sola viga suelta es solo una alerta. La edición avisa
	solo de las vigas que se sueltan por el cambio. \\
\end{longtable}

\begin{nota}
En \ruta{SismoE030Test.\allowbreak{}tablasDeLaNormaCoincidenConValoresConocidos}, las comprobaciones de $T_P = 0{,}6$ y $T_L = 2{,}0$
para S2 sin $V_{s30}$ quedaron por error dentro de un comentario de línea, por lo que hoy no se ejecutan. Los valores sí
se comprueban indirectamente en los casos de $C(1{,}0)$ y $C(3{,}0)$.
\end{nota}

% ------------------------------------------------------------------------------
\subsection{Supuestos y límites}
\label{sec:sis-limites}

\begin{itemize}
	\item \textbf{Sistema:} solo pórticos de concreto armado. Los otros sistemas se pueden elegir para ver su $R_0$, pero
	      no se calculan. Los muros del modelo no aportan rigidez ni cortante.
	\item \textbf{Apoyos:} base empotrada, sin zapatas ni interacción suelo-estructura.
	\item \textbf{Elementos:} barras de Euler--Bernoulli de eje a eje, sin zonas rígidas en los nudos ni deformación por
	      cortante; secciones rectangulares brutas (o reducidas por la fracción de rigidez, con aviso).
	\item \textbf{Diafragma:} rígido en cada piso; las losas no se modelan como elementos, solo como peso.
	\item \textbf{Pisos:} se definen por la cara superior de las vigas con una tolerancia de 5~cm; los niveles muy cercanos
	      se agrupan y las columnas que nacen fuera de un piso se consideran empotradas.
	\item \textbf{Pesos:} 24~kN/m$^3$ para el concreto, sin descontar intersecciones; las cargas por m$^2$ las fija el
	      usuario (las sugeridas pueden no estar verificadas).
	\item \textbf{Análisis:} lineal-elástico; amortiguamiento de 5~\%; excentricidad accidental de 0,05 de la dimensión en
	      planta del rectángulo envolvente; combinación direccional solo en 2026.
	\item \textbf{Verificaciones:} solo deriva de concreto armado (0,007), torsión, cortante mínimo, aplicabilidad, Tabla 9,
	      aislamiento y junta. No se verifican la resistencia ni el diseño de vigas, columnas o nudos, ni las demás
	      irregularidades: $I_a$ e $I_p$ los elige el usuario.
	\item \textbf{Tiempo-historia:} verificación complementaria lineal con registros escalados a $R = 1$; no sustituye al
	      análisis estático ni al espectral.
	\item \textbf{Norma:} los textos y artículos son los que el código tiene escritos; si su proyecto tiene un trámite
	      iniciado con la versión anterior, verifique cuál rige.
\end{itemize}

\begin{advertencia}[Módulo preliminar, no validado]
El análisis sísmico de ConcretoLab Desktop es una herramienta \textbf{educativa y de anteproyecto}. Sus pruebas
automáticas comprueban la programación contra valores conocidos, no constituyen una validación profesional del método.
\textbf{No use estos resultados para diseñar ni para construir obras}, ni los presente como memoria de cálculo ante una
entidad. Todo edificio real debe ser analizado y diseñado por un ingeniero civil colegiado, con un programa de análisis
estructural reconocido y con los criterios completos de la norma vigente.
\end{advertencia}
